% !TeX root = ../../Book4.tex
% 纲要标题：### D.3 短正合列产生的三角
% 纲要细目（写作范围）：
% ---
% #### D.3.1 推出构造与连接态射
% #### D.3.2 三角结构及映射锥

\section{短正合列产生的三角}
\label{b4:sec:D03}

把投射内射对象变成零以后，短正合列的中项和两端之间仍有联系。例如双数代数的列 \(0\to k\to A\to k\to0\) 不能只剩下两个零箭头：它的两个端项由余合冲操作联系起来。下面用推出构造保留这项连接资料，并与复形的映射锥比较。

\subsection{推出构造与连接态射}
\label{b4:subsec:D03:01}

固定 Frobenius 范畴中的序列 \(0\to X\xrightarrow{i_X}I_X\xrightarrow{q_X}\Sigma X\to0\)。对任意 \(f:X\to Y\)，沿 \(f\) 推出可容许单态射 \(i_X\)，得到
\begin{equation}\label{b4:eq:D-standard-pushout}
\begin{tikzcd}[column sep=large,row sep=large]
X\arrow[r,"i_X"]\arrow[d,"f"']&I_X\arrow[r,"q_X"]\arrow[d,"a_f"]&\Sigma X\arrow[d,equal]\\
Y\arrow[r,"b_f"']&C_f\arrow[r,"c_f"']&\Sigma X.
\end{tikzcd}
\end{equation}
下行是可容许短正合列。推出也给出可容许短正合列
\begin{equation}\label{b4:eq:D-graph-inflation}
0\longrightarrow X\xrightarrow{(f,-i_X)}Y\oplus I_X
\xrightarrow{(b_f,a_f)}C_f\longrightarrow0.
\end{equation}
这里第一个映射是可容许单态射：它由分裂包含 \(X\to X\oplus Y\)、变换 \((x,y)\mapsto(x,y+fx)\)、可容许单态射 \((-i_X)\oplus1_Y\) 及交换两项的同构复合而成。推出的泛性质给出第二个映射的余核性质，核与可容许性由正合结构的推出性质保证；其一般表述见 Bühler\cite[Proposition 2.12]{Buehler2010ExactCategories}。

\begin{definition}\label{b4:def:D-stable-standard-triangle}
设 \(\mathcal E\) 为局部小 Frobenius 范畴，固定余合冲序列并由式\eqref{b4:eq:D-standard-pushout}构造 \(C_f\)。在 \(\underline{\mathcal E}\) 中，把
\[
X\xrightarrow{[f]}Y\xrightarrow{[b_f]}C_f\xrightarrow{[c_f]}\Sigma X
\]
称为标准三角。与某个标准三角同构的三角称为好三角。
\end{definition}

若 \(f'\) 与 \(f\) 之差通过 \(P\in\mathcal P\) 分解为 \(vu\)，可把 \(u:X\to P\) 延拓为 \(t:I_X\to P\)。推出的泛性质给出 \(C_f\to C_{f'}\)：在 \(Y\) 上用恒等，在 \(I_X\) 上用 \(a_{f'}-b_{f'}vt\)。反向把负号改为正号，两复合在 \(Y,I_X\) 上都是恒等，因而互逆；它们也保持到 \(\Sigma X\) 的映射。因此标准三角的同构类型只依赖稳定态射。

\begin{proposition}\label{b4:prop:D-conflation-triangle}
设 \(\mathcal E\) 为局部小 Frobenius 范畴，固定余合冲序列。给定可容许短正合列
\[
0\longrightarrow X\xrightarrow{u}Y\xrightarrow{v}Z\longrightarrow0,
\]
把 \(i_X:X\to I_X\) 沿 \(u\) 延拓为 \(a:Y\to I_X\)，并由余核定义 \(\delta:Z\to\Sigma X\)，满足
\begin{equation}\label{b4:eq:D-connecting-sign}
\delta v=-q_Xa.
\end{equation}
则 \([\delta]\) 不依赖延拓的选择，而且
\[
X\xrightarrow{[u]}Y\xrightarrow{[v]}Z\xrightarrow{[\delta]}\Sigma X
\]
是好三角。反过来，每个好三角都同构于某个可容许短正合列按此方式给出的三角。
\end{proposition}
\begin{proof}
由于 \(q_Xau=q_Xi_X=0\)，\(\delta\) 存在且唯一。两个延拓之差经 \(v\) 写成 \(dv\)，其中 \(d:Z\to I_X\)，所以两个连接态射之差为 \(-q_Xd\)，通过 \(I_X\) 分解。

考虑 \(u\) 的推出 \(C_u\)。把原短正合列沿 \(i_X\) 推出，还得到可容许短正合列
\[
0\longrightarrow I_X\xrightarrow{a_u}C_u\xrightarrow{r}Z\longrightarrow0,
\qquad rb_u=v.
\]
推出泛性质由 \(1_{I_X}\) 和 \(a:Y\to I_X\) 给出 \(s:C_u\to I_X\)，满足 \(sa_u=1\)、\(sb_u=a\)。所以此列分裂，\(C_u\cong I_X\oplus Z\)，\([r]\) 在稳定范畴中可逆。在推出的两个来源上分别检查，有
\[
c_u=q_Xs+\delta r:
\quad q_Xsa_u+\delta ra_u=q_X,
\quad q_Xsb_u+\delta rb_u=q_Xa+\delta v=0.
\]
因此 \([c_u]=[\delta]\,[r]\)，而 \([r][b_u]=[v]\)。这给出所写三角与标准三角的同构。

反方向，式\eqref{b4:eq:D-graph-inflation}是可容许短正合列。以 \(-\operatorname{pr}_{I_X}:Y\oplus I_X\to I_X\) 延拓 \(i_X\)，式\eqref{b4:eq:D-connecting-sign}在其余核上给出的连接态射就是 \(c_f\)。在稳定范畴中 \(Y\oplus I_X\cong Y\)，第一个映射成为 \([f]\)，第二个成为 \([b_f]\)，所以该短正合列给出标准三角。
\end{proof}

式\eqref{b4:eq:D-connecting-sign}中的负号配合本卷映射锥的末箭头约定。连接态射的稳定类记录扩张，而不是从原范畴任选一个 \(Z\to\Sigma X\) 的映射。

\subsection{三角结构及映射锥}
\label{b4:subsec:D03:02}

\begin{theorem}[Happel]\label{b4:thm:D-happel-triangulation}
设 \(\mathcal E\) 为局部小 Frobenius 范畴，固定式\eqref{b4:eq:D-syzygy-sequences}中的余合冲序列。以 \(\Sigma\) 为移位、以\cref{b4:def:D-stable-standard-triangle}中的好三角为指定三角，\(\underline{\mathcal E}\) 成为三角范畴。更换余合冲序列后，由比较映射给出的恒等对象函子连同 \(\Sigma\) 的比较自然同构是三角等价。
\end{theorem}

这一结果把正合范畴中的推出与本卷\cref{b4:def:triangulated-category}的三角公理联系起来。定理及复形例见 Bühler\cite[Section 13.4, Theorem 13.7 and Example 13.8]{Buehler2010ExactCategories}；该处的定理追溯到 Happel 专著第一章定理 2.6（第 16 页）的证明。

\begin{proofsketch}
\cref{b4:thm:D-stable-syzygy-equivalence}已经给出移位自等价。每个稳定态射由某个 \(f\) 表示，推出给出补全它的标准三角；\(1_X\) 的推出中 \(C_{1_X}\cong I_X\)，得到 \(X\xrightarrow{1}X\to0\to\Sigma X\)。若比较方块只在稳定范畴中交换，其误差通过一个投射内射对象分解；内射延拓使误差可以吸收到推出映射中，故可以构造三角之间的比较态射。旋转则通过把一条可容许短正合列与 \(Y\to I_Y\) 的序列共同推出，消去新增的投射内射项得到。

八面体所需的对象可以具体构造。对可复合的 \(X\xrightarrow{f}Y\xrightarrow{g}Z\)，令
\[
u=(f,-i_X):X\longrightarrow Y\oplus I_X,
\qquad
v(y,a)=(gy,-i_Yy,a):Y\oplus I_X\longrightarrow Z\oplus I_Y\oplus I_X.
\]
两者均为可容许单态射，\(v\) 是 \((g,-i_Y)\) 与 \(1_{I_X}\) 的直和，因此复合 \(vu\) 也可容许。令 \(T=\operatorname{coker}(vu)\)。对这两个可容许单态射及其复合应用正合范畴的 \(3\times3\) 引理，得到可容许短正合列
\[
0\longrightarrow\operatorname{coker}u\longrightarrow T
\longrightarrow\operatorname{coker}v\longrightarrow0.
\]
这一逐次商结论及 \(3\times3\) 引理分别见 Bühler\cite[Lemma 3.5; Corollary 3.6]{Buehler2010ExactCategories}。这里两端分别是 \(C_f\) 与 \(C_g\)。

为比较 \(T\) 与固定选择的 \(C_{gf}\)，利用 \(I_Y\) 的内射性，把 \(i_Yf:X\to I_Y\) 沿 \(i_X\) 延拓为 \(t:I_X\to I_Y\)，满足 \(ti_X=i_Yf\)。于是
\[
(i_Yf,i_X)=(t,1_{I_X})i_X:X\longrightarrow I_Y\oplus I_X.
\]
其中 \((t,1_{I_X})\) 与标准直和包含同构，所以这也是可容许单态射。复合 \(vu\) 的三个分量为 \((gf,-i_Yf,-i_X)\)。目标上的自同构
\[
\alpha(z,b,a)=(z,b-ta,a),
\qquad
\alpha^{-1}(z,b,a)=(z,b+ta,a)
\]
把它变为 \((gf,0,-i_X)\)，从而
\[
T\cong\operatorname{coker}(gf,0,-i_X)
\cong C_{gf}\oplus I_Y.
\]
因此 \(T\) 与 \(C_{gf}\) 在稳定范畴中同构。上面的可容许短正合列由\cref{b4:prop:D-conflation-triangle}给出八面体的第四个三角
\[
C_f\longrightarrow C_{gf}\longrightarrow C_g\longrightarrow\Sigma C_f.
\]
原来的两个包含与余核图同时提供其余三角间的映射。完整证明还须逐个核对旋转的负号、这些映射的交换性及更换嵌入后的相容性；这些公理核验采用所引 Happel 定理的证明，而以上构造说明八面体中的对象和短正合列从何而来。
\end{proofsketch}

\begin{proposition}\label{b4:prop:D-complex-triangle-identification}
设 \(\mathcal A\) 为局部小加性范畴，\(\operatorname{Ch}(\mathcal A)\) 配备逐项分裂正合结构，并取 \(I_X=\operatorname{Cone}(1_X)\)。则\cref{b4:thm:D-stable-complex-homotopy}中的识别
\[
\underline{\operatorname{Ch}(\mathcal A)}\cong K(\mathcal A)
\]
把 Frobenius 范畴的好三角与映射锥的好三角完全对应。
\end{proposition}
\begin{proof}
对 \(f:X\to Y\)，推出可以取 \(C_f=\operatorname{Cone}(f)\)。其两张结构映射逐次为
\[
Y\longrightarrow\operatorname{Cone}(f),\quad y\longmapsto(y,0),
\qquad
I_X\longrightarrow\operatorname{Cone}(f),\quad (x,x')\longmapsto(fx,x').
\]
它们在 \(X\) 上相同。给定从 \(Y,I_X\) 到复形 \(W\) 的相容映射 \(b,a\)，唯一的逐项推出映射为
\[
\operatorname{Cone}(f)^n\longrightarrow W^n,
\qquad (y,x')\longmapsto b^n(y)+a^n(0,x').
\]
在第二个分量上，其复形条件由 \(a\) 为复形映射及 \(a^n(x,0)=b^n f^n(x)\) 保证；第一个分量由 \(b\) 为复形映射保证。因此这确实是复形范畴中的推出。到 \(\Sigma X=X[1]\) 的映射为 \((y,x')\mapsto x'\)。所得标准三角正是本卷的映射锥三角。两类好三角均由这些标准三角的同构闭包定义，故完全相同。
\end{proof}

对逐项分裂短正合列，若选坐标 \(Y^n=X^n\oplus Z^n\)，写
\[
\mathrm{d}_Y^n=
\begin{pmatrix}\mathrm{d}_X^n&t^n\\0&\mathrm{d}_Z^n\end{pmatrix},
\]
则 \(a^n(x,z)=(x,t^nz):Y^n\to I_X^n\) 是延拓 \(i_X\) 的复形映射。式\eqref{b4:eq:D-connecting-sign}给出连接态射 \(\delta=-t:Z\to X[1]\)。它正是选分裂后由微分的非对角部分得到的映射，换分裂只改变其同伦类代表。

\ProseParagraphBreak
在双数代数中，取 \(i_k(1)=\varepsilon\)、\(I_k=A\) 及 \(\Sigma k=k\)。对式\eqref{b4:eq:D-dual-number-extension}可用 \(a=1_A\)，于是 \(\delta=-1_k\)。稳定范畴中的好三角为
\[
k\longrightarrow0\longrightarrow k\xrightarrow{-1_k}\Sigma k.
\]
因此中项变成零以后，非分裂扩张仍通过非零连接态射留下信息。稳定范畴既能表现模的扩张，也能表现复形的同伦；共同的原因是正合结构中的投射与内射对象可以同时消去。
